\subsection{Taxa de Distração}
Mede a proporção do tempo total de sessão que os alunos passaram em sites classificados como distractores em relação ao tempo total de participação na aula. Esta métrica permite quantificar o nível de foco dos estudantes durante as sessões de aprendizagem. O cálculo é feito conforme a Equação~\ref{eq:student_disturb_ratio} a seguir:
\begin{equation}\label{eq:student_disturb_ratio}
D_i = \frac{T_{d,i}}{T_{s,i}}
\end{equation}

Para um aluno $i$, temos $T_{d,i}$, o tempo total fora da guia de aula, enquanto $T_{s,i}$ diz respeito ao tempo total de sessão (em minutos). 

De posse dos dados individuais do aluno $i$, pode-se calcular a taxa de duranção média correspondente à turma, definida pela Equação~\ref{eq:class_mean_disturb_ratio}.

\begin{equation}\label{eq:class_mean_disturb_ratio}
\bar{D} = \frac{1}{n} \sum_{i=1}^{n} D_i
\end{equation}

% \lstinputlisting[
%     language=python,
%     caption={Cálculo da Taxa de Distração},
%     firstline=7,
%     lastline=65,
%     label={lst:calculate_attendance_ratio}
% ]{scripts/metricsCalc/disturb.py}

Primeiramente é feito o agrupamento dos \textit{traces} por usuário, seguido da sua ordenação por \textit{timestamp}. Feito isso, para cada intervalo entre \textit{traces} consecutivos é caulculado intervalo etre os registros (em minutos), que é dividido pelo tempo total da sessão. Traces considerados distração, isto é, que não a aula, são adicionados à pilha para cálculo do tempo total de \textit{disturbance}. Ao fim, é feito o cálculo da proporção de distração para cada aulo, seguido da emissão de média aritmética das taxas individuais (aluno).

% \textbf{Funções auxiliares}:
% \begin{itemize}
%     \item \texttt{get\_time\_diff\_minutes()}: Calcula diferença temporal entre \textit{timestamps}
%     \item \texttt{is\_distraction()}: Classifica se um \textit{trace} representa atividade em site distrator
% \end{itemize}

\subsection{Frequência de Distração}
Esta métrica mede quantas vezes por hora os alunos alternam da aula para sites, que não a guia de aula, complementando a \textbf{Taxa de Distração}, capturando a dinâmica temporal das transições de atenção, ao invés de apenas a proporção total de tempo distraído. Seu cálculo está disposto na Equação~\ref{eq:disturbance_ration} abaixo:

\begin{equation}\label{eq:disturbance_ration}
\text{Frequência de Distração} = \frac{\text{Número de Transições para Distração}}{\text{Tempo Total de Sessão (horas)}}
\end{equation}

% \lstinputlisting[
%     language=python,
%     caption={Cálculo da Taxa de Distração},
%     firstline=68,
%     lastline=129,
%     label={lst:calculate_distraction_frequency}
% ]{scripts/metricsCalc/disturb.py}

O processo de cálculo consiste em agrupar e ordenar cronologicamente os \textit{traces} por aluno, identificando transições entre os estados “em aula” e “distração” em registros consecutivos. Em seguida, calcula-se o tempo total de sessão por meio do somatório dos intervalos entre registros válidos. As transições detectadas são então normalizadas em frequência por hora, permitindo comparar sessões de durações distintas. Por fim, obtém-se um valor agregado da turma a partir da média das frequências individuais dos alunos.

% \textbf{Funções auxiliares}:
% \begin{itemize}
%     \item \textbf{is\_lecture\_tab()}: Determina se o \textit{trace} representa atividade na aba da aula
%     \item \textbf{is\_distraction()}: Classifica se o \textit{trace} representa atividade em site distrator
%     \item \textbf{get\_time\_diff\_minutes()}: Calcula diferença temporal entre \textbf{timestamps}
% \end{itemize}